\section{Introduction}
\label{sec:introduction}

Adversarial and distribution-shift failures remain a practical barrier for deploying image classifiers in high-stakes settings. Deep ensembles improve uncertainty quality, but training many independent models is expensive \citep{lakshminarayanan2017deep}. Neural thicket style methods suggest a cheaper alternative: sample diverse experts around one pretrained solution in weight space \citep{gan2026neuralthickets,mehrtash2020pep}.

\para{{\bf why is this still unresolved?}} Existing studies show that diversity can help uncertainty estimation, but the link between perturbation-generated diversity and robustness is not settled under a controlled adversarial + OOD protocol \citep{xia2022usefulness,wei2020xensemble}. In particular, it is unclear whether structured perturbation directions are better than random directions when ensemble size and perturbation scale are fixed.

We examine this question using one CIFAR-10 ResNet-18 base model and three perturbation strategies: random, orthogonal, and adversarial-direction perturbations. We evaluate clean accuracy, calibration, PGD robustness at $\epsilon\in\{2/255,4/255,8/255\}$, and OOD detection on SVHN. This design isolates the effect of direction strategy while keeping architecture, data, and expert count constant.

Our key quantitative result is narrow but clear. Random perturbation ensembling improves PGD@8/255 from $0.05\%$ to $0.50\%$ (absolute $+0.45$ points, 95\% CI $[+0.15,+0.80]$), while reducing clean accuracy from $88.28\%$ to $86.99\%$. Orthogonal directions do not significantly improve over random at PGD@8/255 ($p=0.599$), and adversarial-direction perturbations degrade calibration and robustness.

Our contributions are:
\begin{itemize}[leftmargin=*,itemsep=0pt,topsep=0pt]
    \item We conduct a controlled robustness benchmark for perturbation-based thicket ensembles across clean, adversarial, OOD, and calibration metrics.
    \item We compare random, orthogonal, and adversarial-direction perturbation strategies under matched ensemble budget ($M=6$) and perturbation scale ($\sigma=0.015$).
    \item We provide bootstrap confidence intervals and hypothesis tests for primary robustness claims, including negative findings.
    \item We show that larger disagreement alone does not predict better robustness or OOD separation in this setup.
\end{itemize}

\para{Paper organization.} \secref{sec:related_work} reviews prior work. \secref{sec:methodology} describes datasets, models, and protocol. \secref{sec:results} presents quantitative results and statistical tests. \secref{sec:discussion} discusses interpretation and limitations, and \secref{sec:conclusion} concludes.
